\medskip

Un étudiant mange tous les jours au restaurant universitaire. Ce restaurant propose des plats végétariens et des plats non végétariens.

\begin{itemize}[label=$\bullet~~$]
\item Lorsqu'un jour donné l'étudiant a choisi un plat végétarien, la probabilité qu'il choisisse un plat végétarien le lendemain est $0,9$.
\item Lorsqu'un jour donné l'étudiant a choisi un plat non végétarien, la probabilité qu'il choisisse un plat végétarien le lendemain est $0,7$.
\end{itemize}

\medskip

Pour tout entier naturel $n$, on note $V_n$ l'évènement \og l'étudiant a choisi un plat végétarien le $n^{\text{e}}$ jour \fg{} et $p_n$ la probabilité de $V_n$.

Le jour de la rentrée, l'étudiant a choisi le plat végétarien. On a donc $p_1 = 1$.

\medskip

\begin{enumerate}
\item
	\begin{enumerate}
		\item Indiquer la valeur de $p_2$.
		\item Montrer que $p_3 = 0,88$. On pourra s'aider d'un arbre pondéré.
		\item Sachant que le $3\up{\text{e}}$ jour l'étudiant a choisi un plat végétarien, quelle est la probabilité qu’il ait choisi un plat non végétarien le jour précédent ?

On arrondira le résultat à $10^{-2}$.
	\end{enumerate}
\item Recopier et compléter l'arbre pondéré ci-dessous:

\begin{center}
\pstree[treemode=R,levelsep=3cm]{\TR{}}
{\pstree{\TR{$V_n~$}\taput{$p_n$}}
	{\TR{$V_{n+1}$}
	\TR{$\overline{V_{n+1}}$} 
	}
\pstree{\TR{$\overline{V_n}~$}}
	{\TR{$V_{n+1}$}
	\TR{$\overline{V_{n+1}}$} 
	}
}
\end{center}
\item Justifier que, pour tout entier naturel $n \geqslant 1,\:\: p_{n+1} = 0,2p_n + 0,7$.
\item On souhaite disposer de la liste des premiers termes de la suite $(p_n)$ pour $n \geqslant 1$.

Pour cela, on utilise une fonction appelée \texttt{repas} programmée en langage Python dont on propose trois versions, indiquées ci-dessous.
\end{enumerate}
\hspace{-0.5cm}\vspace{-0.5cm}

\begin{center}
{\footnotesize
\begin{ttfamily}
\begin{tabularx}{\linewidth}{*{3}{>{\centering \arraybackslash}X}}
\multicolumn{1}{c}{Programme 1}&\multicolumn{1}{c}{Programme 2}&\multicolumn{1}{c}{Programme 3}\\ 
\begin{tabular}{l| l|}\cline{2-2}
1&def repas(n):\\
2&~~p=1\\
3&~~L=[p]\\
4&~~for k in range(1,n):\\
5&~~~~p = 0.2*p+0.7\\
6&~~~~L. append(p)\\
7&~~~return(L)\\ \cline{2-2}
\end{tabular}&
\begin{tabular}{l| l|}\cline{2-2}
1&def repas(n):\\
2&~~p=1\\
3&~~L=[p]\\
4&~~for k in range(1,n+1):\\
5&~~~~p = 0.2*p+0.7\\
6&~~~~L. append(p)\\
7&~~return(L)\\ \cline{2-2}
\end{tabular}&
\begin{tabular}{l| l|}\cline{2-2}
1&def repas(n):\\
2&~~p=1\\
3&~~L=[p]\\
4&~~for k in range(1,n):\\
5&~~~~p = 0.2*p+0.7\\
6&~~~~L.append(p+1)\\
7&~~return(L)\\ \cline{2-2}
\end{tabular}\\
\end{tabularx}
\end{ttfamily}
}
\end{center}

\begin{enumerate}[start=4]
	\item
	\begin{enumerate}
		\item Lequel de ces programmes permet d’afficher les $n$ premiers termes de la suite $(p_n)$ ? Aucune justification n'est attendue.
		\item Avec le programme choisi à la question \textbf{a.} donner le résultat affiché pour $n = 5$.
	\end{enumerate}
\item Démontrer par récurrence que, pour tout naturel $n \geqslant 1,\:\: p_n = 0,125 \times 0,2^{n - 1} + 0,875$.
\item En déduire la limite de la suite $(p_n)$.
\end{enumerate}

\bigskip
